%%%%%%%%%%%%%%%%%%%%%%%%%%%%%%%%%%%%%%%%%%%%%%%%%%%%%%%%%%%%%%%%%%%%%%%%
%% Art. 27, II a IV — a seção mais decisiva para o exame:              %%
%%   II  descrever o estado da técnica útil à compreensão, à busca e ao %%
%%       exame do pedido, citando sempre que possível os documentos que %%
%%       o reflitam, e destacando os problemas técnicos existentes;      %%
%%   III descrever, de forma clara, concisa e precisa, a solução        %%
%%       proposta e as vantagens em relação ao estado da técnica;        %%
%%   IV  ressaltar nitidamente a novidade e evidenciar o efeito técnico %%
%%       alcançado (invenção) ou a melhoria funcional (modelo).          %%
%%                                                                     %%
%% Como citar o estado da técnica: em prosa, dentro dos parágrafos      %%
%% numerados, identificando o documento pelo número de publicação (ex.: %%
%% BR 10 2015 000000 0, US 2018/0123456 A1, EP 3 000 000 B1). NÃO há    %%
%% lista de referências no pedido de patente — nada de ABNT aqui. As    %%
%% diretrizes de exame (Resolução INPI/PR nº 124/2013, itens 2.40 a     %%
%% 2.46) tratam de como os documentos citados são considerados.         %%
%%                                                                     %%
%% Ponto de atenção do art. 32 da LPI, regulado pela Resolução          %%
%% INPI/PR nº 93/2013: informação essencial ao exame e à                %%
%% patenteabilidade NÃO pode ser inserida depois do requerimento de     %%
%% exame. O que não estiver aqui até lá pode custar o indeferimento,    %%
%% mesmo que o objeto seja novo e inventivo.                            %%
%%%%%%%%%%%%%%%%%%%%%%%%%%%%%%%%%%%%%%%%%%%%%%%%%%%%%%%%%%%%%%%%%%%%%%%%

\secaoinpi{Fundamentos \danatureza}

\pnum São conhecidos do estado da técnica dispositivos de acionamento de
válvulas de irrigação que empregam solenoides de acionamento direto, nos quais
a força magnética atua sobre a haste de vedação ao longo de todo o curso de
abertura. O documento BR 10 2015 000000 0 descreve um conjunto desse tipo, no
qual o solenoide permanece energizado durante todo o intervalo de irrigação.

\pnum Também são conhecidos dispositivos de acionamento hidráulico assistido,
como o descrito no documento US 2018/0123456 A1, em que a pressão da própria
linha auxilia o deslocamento do elemento de vedação, reduzindo a força exigida
do atuador.

\pnum Os dispositivos de acionamento direto do estado da técnica apresentam
consumo elevado de energia, uma vez que exigem corrente contínua de manutenção
durante todo o período de abertura da válvula, o que limita seu emprego em
sistemas alimentados por painéis fotovoltaicos de pequeno porte. Os
dispositivos de acionamento hidráulico assistido, por sua vez, dependem de uma
pressão mínima de linha para operar, e apresentam falha de fechamento quando a
pressão cai abaixo desse limite.

\pnum O problema técnico não resolvido pelo estado da técnica consiste,
portanto, em acionar a válvula com baixo consumo de energia sem tornar o
fechamento dependente da pressão da linha.

\ifinvencao
\pnum A presente invenção resolve esse problema por meio de um dispositivo de
acionamento biestável, no qual o atuador eletromecânico atua apenas durante a
transição entre os estados de abertura e de fechamento, e a mola de retorno
mantém o estado alcançado sem consumo de energia. O efeito técnico alcançado é
a redução do consumo de energia por ciclo de acionamento, com fechamento
garantido por meio mecânico, independentemente da pressão da linha.

\pnum A novidade da invenção reside na associação entre o assento de vedação de
dupla face e a mola de retorno pré-carregada, que juntos estabelecem dois
estados estáveis de operação, dispensando a corrente de manutenção exigida
pelos dispositivos de acionamento direto conhecidos.
\else
\pnum O presente modelo de utilidade resolve esse problema por meio de nova
disposição construtiva do conjunto de acionamento, na qual a haste de
acionamento recebe um ressalto de encosto que atua sobre a mola de retorno
pré-carregada, de modo que o conjunto passa a apresentar dois estados estáveis
de operação. A melhoria funcional alcançada é a redução do consumo de energia
por ciclo de acionamento, com fechamento garantido por meio mecânico,
independentemente da pressão da linha.

\pnum A nova forma reside no ressalto de encosto da haste de acionamento e em
seu posicionamento em relação ao assento de vedação de dupla face, disposição
que não se encontra nos objetos conhecidos do estado da técnica.
\fi
